% ===========================================================================
%  Generated by scripts/nb2tex.py from
%    Adv Lesson 01 Exercises.nb
%  Do not edit this file: regenerate it. Output cells are dropped by design --
%  the submission format for this course is a printed notebook with the output
%  cleared (README, "Submission format").
% ===========================================================================
\documentclass[11pt]{article}
\usepackage{coursemacros}
\usepackage{alltt}

\setlength{\parindent}{0pt}
\setlength{\parskip}{0.6em}

% The notebook numbers its own sections ("1: Front-end tour"). Numbering them
% again gives "1  1: Front-end tour", so the sections here are unnumbered and
% the notebook's numbers are left to speak for themselves.
\pagestyle{fancy}
\fancyhf{}
\fancyhead[L]{Advanced Lesson 1 - Exercises: The Mathematica Interface}
\fancyhead[R]{Mikey Ferguson}
\fancyfoot[C]{\thepage}
\renewcommand{\headrulewidth}{0.4pt}

\begin{document}

\begin{center}
  {\Large\bfseries Advanced Lesson 1 - Exercises: The Mathematica Interface}\\[0.4em]
  {\normalsize MATH 7013 - Advanced Mathematical Modeling}\\[0.3em]

  {\normalsize Mikey Ferguson}\\[0.2em]
  {\small \today}
\end{center}
\vspace{0.5em}

\section*{1: Front-end tour - toolbar, magnification, Documentation Center}

This exercise is entirely a walk through the notebook front end: show the
toolbar, raise the magnification, open the Documentation Center, look up
Integrate, and read the seven linked tutorials (Your First Mathematica
Calculations, Arithmetic, Defining Variables, Defining Functions, Some
Mathematical Functions, Symbolic Calculations, Numerical Calculations).
The professor's own instruction is ``Don't worry about an answer--just do it!''
There is nothing to evaluate and nothing to submit.

One thing worth doing while the Documentation Center is open: the Basic Math
Input palette (Palettes menu) supplies the typeset forms this lesson keeps
asking for - the superscript template, the partial-derivative template, the
integral-with-differential template, and the Pi, E and I characters. Every
one of those has a keyboard equivalent, and the exercises below are written
in ordinary linear syntax, which evaluates identically.

\section*{2: The numerical function N}

\begin{quote}\ttfamily\small\begin{alltt}
N[1/13]
\end{alltt}\end{quote}


\begin{quote}\ttfamily\small\begin{alltt}
N[1/13, 50]
\end{alltt}\end{quote}


N with one argument returns a machine-precision real: about sixteen digits,
carried in hardware. N with a second argument asks for that many digits of
precision and Mathematica switches to software arbitrary-precision
arithmetic to supply them.

\section*{3: Assignment, and what ``?'' reports}

\begin{quote}\ttfamily\small\begin{alltt}
Clear[a, A];
a A
\end{alltt}\end{quote}


\begin{quote}\ttfamily\small\begin{alltt}
?a
\end{alltt}\end{quote}


\begin{quote}\ttfamily\small\begin{alltt}
?A
\end{alltt}\end{quote}


\begin{quote}\ttfamily\small\begin{alltt}
a=5
\end{alltt}\end{quote}


\begin{quote}\ttfamily\small\begin{alltt}
a A
\end{alltt}\end{quote}


\begin{quote}\ttfamily\small\begin{alltt}
?a
\end{alltt}\end{quote}


\begin{quote}\ttfamily\small\begin{alltt}
?A
\end{alltt}\end{quote}


\begin{quote}\ttfamily\small\begin{alltt}
A=6
\end{alltt}\end{quote}


\begin{quote}\ttfamily\small\begin{alltt}
a A
\end{alltt}\end{quote}


\begin{quote}\ttfamily\small\begin{alltt}
?a
\end{alltt}\end{quote}


\begin{quote}\ttfamily\small\begin{alltt}
?A
\end{alltt}\end{quote}


\begin{quote}\ttfamily\small\begin{alltt}
A=7
\end{alltt}\end{quote}


\begin{quote}\ttfamily\small\begin{alltt}
a A
\end{alltt}\end{quote}


\begin{quote}\ttfamily\small\begin{alltt}
?a
\end{alltt}\end{quote}


\begin{quote}\ttfamily\small\begin{alltt}
?A
\end{alltt}\end{quote}


Results. With nothing assigned, a*A stays as the symbolic product a A, and
``?'' reports each symbol with no value attached. After a = 5 the product
becomes 5 A: Mathematica substitutes every value it has and leaves the rest
symbolic. After A = 6 the product is 30, and after A = 7 it is 35 - the new
assignment to A simply overwrites the old one. Note that a and A are
different symbols; Mathematica is case sensitive, which is exactly why the
convention is that built-in names start with a capital and yours do not.

\section*{4: Sin versus sin}

\begin{quote}\ttfamily\small\begin{alltt}
?Sin
\end{alltt}\end{quote}


\begin{quote}\ttfamily\small\begin{alltt}
?sin
\end{alltt}\end{quote}


``?Sin'' reports the built-in: Sin[z] gives the sine of z, with a link into
the Documentation Center. ``?sin'' reports a bare symbol with no value, no
definition and no usage message, because no such built-in exists. Lowercase
sin is just an ordinary undefined symbol - which is precisely why user
functions are given lowercase names: there is no chance of colliding with a
protected built-in.

\section*{5: Symbols as units}

\begin{quote}\ttfamily\small\begin{alltt}
Clear[a, A, feet, sec];
a/A
\end{alltt}\end{quote}


\begin{quote}\ttfamily\small\begin{alltt}
?a
\end{alltt}\end{quote}


\begin{quote}\ttfamily\small\begin{alltt}
?A
\end{alltt}\end{quote}


\begin{quote}\ttfamily\small\begin{alltt}
A=(5 feet)/sec
\end{alltt}\end{quote}


\begin{quote}\ttfamily\small\begin{alltt}
a/A
\end{alltt}\end{quote}


\begin{quote}\ttfamily\small\begin{alltt}
a=10 feet
\end{alltt}\end{quote}


\begin{quote}\ttfamily\small\begin{alltt}
a/A
\end{alltt}\end{quote}


\begin{quote}\ttfamily\small\begin{alltt}
?feet
\end{alltt}\end{quote}


\begin{quote}\ttfamily\small\begin{alltt}
?sec
\end{alltt}\end{quote}


Results. With nothing assigned, a/A is the symbolic quotient. After
A = 5 feet/sec the quotient becomes a sec / (5 feet). After a = 10 feet the
quotient becomes 2 sec: the symbol ``feet'' cancels between numerator and
denominator by ordinary algebra, leaving a pure time. ``?feet'' and ``?sec''
report symbols with no values - they never needed any. That is the whole
trick. Undefined symbols behave as formal algebraic units and cancel
correctly without any units package at all.

\section*{6: Defining a function, differentiating, integrating, solving, plotting}

\begin{quote}\ttfamily\small\begin{alltt}
Clear[x, f];
f[x\_]=Sin[x];
\end{alltt}\end{quote}


\begin{quote}\ttfamily\small\begin{alltt}
\{f[Pi], f[Pi/2], f[6.]\}
\end{alltt}\end{quote}


f[Pi] returns exactly 0 and f[Pi/2] returns exactly 1, because Pi is an
exact symbolic constant and Mathematica knows those two values. f[6.0]
returns an approximate real, because the argument was approximate: an
inexact input forces an inexact output.

\begin{quote}\ttfamily\small\begin{alltt}
?f
\end{alltt}\end{quote}


\begin{quote}\ttfamily\small\begin{alltt}
D[f[x], x]
\end{alltt}\end{quote}


\begin{quote}\ttfamily\small\begin{alltt}
f'[x]
\end{alltt}\end{quote}


\begin{quote}\ttfamily\small\begin{alltt}
Integrate[f[x], x]
\end{alltt}\end{quote}


\begin{quote}\ttfamily\small\begin{alltt}
Solve[3 f[x]+4==0, x]
\end{alltt}\end{quote}


\begin{quote}\ttfamily\small\begin{alltt}
Plot[f[x], \{x, 0, 2 Pi\}]
\end{alltt}\end{quote}


Both derivative notations return Cos[x]; f'[x] is shorthand for the same
Derivative operation. The integral is -Cos[x], with no constant of
integration - Mathematica returns one antiderivative, not the family.
Solve does not return a single principal value. It returns two complete
solution families, each a ConditionalExpression carrying the side condition
that the arbitrary integer C[1] is an integer: x = -ArcSin[4/3] + 2 Pi C[1]
and x = Pi + ArcSin[4/3] + 2 Pi C[1]. That is the honest answer for a
periodic equation - infinitely many solutions, parameterised. Note also that
4/3 lies outside the range of the real sine, so ArcSin[4/3] is complex and
every one of those solutions is a complex number.

\section*{7: f[x] = x\textasciicircum{}2 versus f[x\_] = E\textasciicircum{}x versus f[2] = Pi}

\begin{quote}\ttfamily\small\begin{alltt}
Clear[a, x, f, y];
f[x]=x\textasciicircum{}2;
f[x\_]=E\textasciicircum{}x;
f[2]=Pi;
\end{alltt}\end{quote}


\begin{quote}\ttfamily\small\begin{alltt}
\{f[x], f[a], f[y], f[2], f[3]\}
\end{alltt}\end{quote}


\begin{quote}\ttfamily\small\begin{alltt}
?f
\end{alltt}\end{quote}


Three separate definitions are now stored, and Mathematica orders them from
most specific to most general: f[2] first, then f[x] (a literal argument,
not a pattern), then f[x\_].
f[x] returns x\textasciicircum{}2 - the literal definition wins for that one symbol.
f[a] and f[y] return E\textasciicircum{}a and E\textasciicircum{}y - no literal definition matches, so the
pattern definition applies.
f[2] returns Pi - the literal definition for 2 wins over the pattern.
f[3] returns E\textasciicircum{}3 - only the pattern matches.
The lesson: f[x] = x\textasciicircum{}2 without the underscore defines the function at the
single symbol x and nowhere else. It is almost never what you want.

\section*{8: Immediate versus delayed assignment with RandomReal}

\begin{quote}\ttfamily\small\begin{alltt}
Clear[f1, f2, x];
\end{alltt}\end{quote}


\begin{quote}\ttfamily\small\begin{alltt}
?RandomReal
\end{alltt}\end{quote}


\begin{quote}\ttfamily\small\begin{alltt}
f1[x\_]=RandomReal[] x;
\{f1[x], f1[x], f1[x]\}
\end{alltt}\end{quote}


\begin{quote}\ttfamily\small\begin{alltt}
Plot[f1[x], \{x, 0, 1\}]
\end{alltt}\end{quote}


\begin{quote}\ttfamily\small\begin{alltt}
?f1
\end{alltt}\end{quote}


\begin{quote}\ttfamily\small\begin{alltt}
f2[x\_]:=RandomReal[] x;
\{f2[x], f2[x], f2[x]\}
\end{alltt}\end{quote}


\begin{quote}\ttfamily\small\begin{alltt}
Plot[f2[x], \{x, 0, 1\}]
\end{alltt}\end{quote}


\begin{quote}\ttfamily\small\begin{alltt}
?f2
\end{alltt}\end{quote}


Results. f1 gives the same coefficient all three times, and its plot is a
straight line through the origin: with immediate assignment the right-hand
side was evaluated once, at definition time, so RandomReal ran once and the
number it produced was frozen into the definition. ``?f1'' shows a definite
numeric coefficient stored in the rule.
f2 gives three different values, and its plot is noise: with delayed
assignment the right-hand side is re-evaluated at every call, so RandomReal
runs afresh for each of the several hundred sample points Plot requests.
``?f2'' shows RandomReal[]*x stored literally, unevaluated.

\section*{9: Derivative, integral, radical equation, and a cubic plot}

\begin{quote}\ttfamily\small\begin{alltt}
Clear[x, n];
D[x\textasciicircum{}n, x]
\end{alltt}\end{quote}


\begin{quote}\ttfamily\small\begin{alltt}
Integrate[x\textasciicircum{}n, x]
\end{alltt}\end{quote}


\begin{quote}\ttfamily\small\begin{alltt}
Solve[x+Sqrt[x]==9, x]
\end{alltt}\end{quote}


\begin{quote}\ttfamily\small\begin{alltt}
Plot[x\textasciicircum{}3-2 x+1, \{x, -2, 2\}]
\end{alltt}\end{quote}


The derivative is n x\textasciicircum{}(n-1) as expected. The integral is x\textasciicircum{}(n+1)/(n+1),
returned without any condition - Mathematica quietly gives the generic
answer and does not stop to warn that n = -1 is excluded. The radical
equation has the single solution x = (19 - Sqrt[37])/2, and Solve discards
the extraneous root that squaring would introduce. The cubic has a local
maximum near x = -0.8 and a local minimum near x = 0.8, with one real root
in view to the left.

\section*{10: Wildcard searches on the name Style}

\begin{quote}\ttfamily\small\begin{alltt}
?Style
\end{alltt}\end{quote}


\begin{quote}\ttfamily\small\begin{alltt}
Names[``Style*'']
\end{alltt}\end{quote}


\begin{quote}\ttfamily\small\begin{alltt}
Length[Names[``Style*'']]
\end{alltt}\end{quote}


\begin{quote}\ttfamily\small\begin{alltt}
Names[``*Style'']
\end{alltt}\end{quote}


\begin{quote}\ttfamily\small\begin{alltt}
Length[Names[``*Style'']]
\end{alltt}\end{quote}


\begin{quote}\ttfamily\small\begin{alltt}
Names[``*Style*'']
\end{alltt}\end{quote}


\begin{quote}\ttfamily\small\begin{alltt}
Length[Names[``*Style*'']]
\end{alltt}\end{quote}


The wildcard ``*'' stands for any run of characters, so ``Style*'' finds names
beginning with Style, ``*Style'' finds names ending in Style, and ``*Style*''
finds names containing it anywhere. Names returns the list; Length counts
it. In this version of the language the counts are 13, 107 and 137. The
exercise quotes ``about 12, about 45, about 65'', which were the counts when
it was written; the built-in name list has grown substantially since, and
the ``ends in Style'' family in particular has more than doubled. The first
count is still essentially right because Style* names are core language
constructs, while the *Style names are mostly plot and graph options, which
is where most of the growth has happened.

\end{document}
